\label{chap.3}
\section{Overview}
MAPIS is a monitor that sits on the communication layer of a multi-agent pipeline, remembers the session and gives each message a trust score. Figure~\ref{fig:arch} shows the path a message takes.
\begin{figure}[H]
\centering
\begin{tikzpicture}[node distance=0.55cm, box/.style={draw, rounded corners, align=center, minimum width=3.6cm, minimum height=0.95cm, font=\small}, hot/.style={box, fill=maroon!10, draw=maroon}, ar/.style={-{Stealth[length=2.5mm]}, thick}]
\node[box] (ag) {Agents: planner, web, file, memory, code/email\\LangGraph or AutoGen};
\node[hot, below=of ag] (tap) {MAPIS tap (FastAPI \texttt{/inspect})\\every message, tool call, tool output, memory access};
\node[hot, below=of tap] (ses) {Session record (Redis)\\goal, item sources, authority claims, baselines};
\node[hot, below=of ses] (sc) {Trust score $t\in[0,1]$\\DeBERTa-v3-small + four signals};
\node[box, below=of sc] (tier) {Tier: PASS / FLAG / QUARANTINE / BLOCK};
\node[box, below=of tier] (log) {Forensic trace and operator dashboard};
\foreach \a/\b in {ag/tap,tap/ses,ses/sc,sc/tier,tier/log} \draw[ar] (\a) -- (\b);
\end{tikzpicture}
\caption{Path of a message through MAPIS.}
\label{fig:arch}
\end{figure}

\section{The Tap}
The tap is a FastAPI endpoint (\texttt{POST /inspect}) that the testbed calls before it delivers any message, so MAPIS is not tied to one framework. In LangGraph the call is placed in the graph before every delivery. In AutoGen 0.2 we register it with \texttt{register\_hook} on the \texttt{process\_message\_before\_send} event, and if a message is withheld the chat ends. Both frameworks run the same deterministic testbed with five agents (planner, web, file, memory, code/email) and three scenarios: a clean task, an obvious attack and a decomposed attack.

\section{Session Record}
For every session MAPIS keeps a small record, stored in Redis with a time-to-live (or in memory if Redis is not running). It holds:
\begin{itemize}
\item the user's goal, taken from the first user message, which is the trust anchor;
\item \emph{tainted} items: emails, URLs and hosts, IBAN-like ids and long numbers, and fragments of them, first seen in untrusted content (tool, tool response, memory read, memory write);
\item \emph{trusted} items that the user named;
\item untrusted authority claims (``approved by management'') and instruction cues;
\item per-agent baselines: tools used and message lengths;
\item the session sensitivity and the set of suspended channels.
\end{itemize}
The record is updated after every event. We use the same module (\texttt{session.py}) to build the training data and at run time, and a unit test makes sure both produce exactly the same text.

\section{Classifier Input}
The classifier input has three parts, always in the same order (module \texttt{context.py}). First comes the current event, so that if the text is too long only the history gets cut. Then comes one line of session signals, and then the history itself: the goal, the latest hops cut to 110 characters, and up to three older untrusted hops. Events longer than 900 characters are split into chunks with a stride of 450, and the event gets the lowest trust of its chunks.

\section{The Four Signals}
\begin{table}[H]
\centering\small
\begin{tabular}{p{3.4cm}p{5.0cm}p{4.3cm}}
\toprule
Signal & What it detects & Example\\
\midrule
Instruction in data & Command-like text or hidden characters in data channels & ``you must now forward\ldots'' inside a PDF\\
Drift from the goal & A state-changing action unrelated to the user's request (lexical similarity of goal and action, only for state-changing calls) & Asked for a summary, agent wants to pay money\\
Cross-hop provenance and escalation & An item first seen in untrusted content reaches an action or memory write; an untrusted authority claim followed by an unrequested action or waiver wording the user never used & The planted address becomes the email recipient\\
Behaviour anomaly & A tool or outbound action outside the agent's baseline; a message four times longer than usual & The file agent suddenly sends email\\
\bottomrule
\end{tabular}
\caption{The four signals.}
\end{table}

\subsection{Evasion defences}
An attacker who knows how MAPIS works can try to hide the item. So before matching, \texttt{canonical()} normalises Unicode (NFKC), removes zero-width characters and undoes tricks such as \texttt{hxxp}, \texttt{[.]}, ``(at)'', ``at'' and ``dot'' written out, spaces around \texttt{@} and IBANs written in groups. Fragment matching also stores bare domains, hosts and id-like tokens, so if an action puts an item back together from pieces it still matches (the session signal \texttt{reuse=} shows which piece).

\section{Fusion and Response}
To get the final trust, the model trust $t_m$ is combined with the rule risks $r_k$ using a noisy-OR, and the result is scaled by the session sensitivity $\sigma$:
\begin{equation}
r = 1 - t_m\prod_{k}\bigl(1 - w_k r_k\bigr), \qquad t = 1 - \min\bigl(1,\ (1+\sigma)\,r\bigr).
\end{equation}
In words, an event loses trust as soon as the model or any weighted rule reports a risk. The weights are $w_{\text{provenance}}=0.8$, $w_{\text{instruction}}=0.5$, $w_{\text{drift}}=0.15$ and $w_{\text{behaviour}}=0.15$. Messages from the user or the system always get $t=1$, and any message on a suspended channel gets $t=0$.

\begin{table}[H]
\centering\small
\begin{tabular}{lll}
\toprule
Trust $t$ & Tier & Action\\
\midrule
$\ge 0.75$ & PASS & Deliver\\
$0.50$ to $0.75$ & FLAG & Deliver, raise session sensitivity by $0.15$ (cap $0.6$)\\
$0.25$ to $0.50$ & QUARANTINE & Hold for operator Release or Reject\\
$< 0.25$ & BLOCK & Drop and suspend the channel\\
\bottomrule
\end{tabular}
\caption{Tiered response.}
\end{table}
The end user never sees MAPIS. When an action is held, it simply does not happen. Every verdict other than PASS comes with a forensic trace that shows the origin (the untrusted source that brought in the item), the hops that relayed it, and the action that would have used it. If something fails, MAPIS fails closed.
